% Generated draft: requires mathematical and source review.
\documentclass[11pt]{article}
\usepackage[margin=25mm]{geometry}
\usepackage{fontspec}
\setmainfont{Latin Modern Roman}
\usepackage{amsmath,amssymb,mathtools,unicode-math}
\setmathfont{Latin Modern Math}
\usepackage{enumitem,needspace,xurl,hyperref}
\setlength{\parindent}{0pt}
\setlength{\parskip}{5pt}
\setlength{\emergencystretch}{3em}
\newcommand{\N}{\mathbb{N}}
\newcommand{\Z}{\mathbb{Z}}
\newcommand{\Q}{\mathbb{Q}}
\newcommand{\R}{\mathbb{R}}
\newcommand{\C}{\mathbb{C}}
\newcommand{\F}{\mathbb{F}}
\newcommand{\A}{\mathbb{A}}
\newcommand{\PP}{\mathbb{P}}
\newcommand{\E}{\mathbb{E}}
\newcommand{\eps}{\varepsilon}
\newcommand{\dd}{\,\mathrm{d}}
\begin{document}
\section*{UnsolvedMath ID: 30001975}
\textbf{Source problem code:} catalogue entry
\subsection{Definitions and mathematical statement}
\textbf{Scope and terminology.} The supplied record refers to Polanyi's concept of tacit knowledge and to explicit mathematical knowledge, but supplies no precise definitions of either concept, no specified historical corpus, and no criterion for explanatory success. These terms therefore remain conceptual terms from the source, rather than formally defined mathematical objects.

\textbf{Source question.} ``Does Polanyi's concept of tacit knowledge help explain the historical development of mathematics, and how does tacit knowledge relate to explicit mathematical knowledge?''

\textbf{Unresolved formulation.} This is a historical, pedagogical, or methodological discussion question. The supplied evidence does not specify a mathematical proposition, domains of quantification, hypotheses, or a conclusion that could be proved or disproved. Consequently, no precise mathematical statement can be recovered without adding substantive assumptions absent from the record. The supplied statement-verification note classifies the formulation as unrecoverable and retains the original question for provenance; it does not promote a separate clean statement. The registry's supplied \emph{open} label should not be interpreted as verification that this is an open mathematical conjecture.
\subsection{Short English statement}
The source asks whether Polanyi's tacit knowledge helps explain mathematical history and how it relates to explicit mathematical knowledge. It supplies a discussion prompt without the definitions or criteria needed for a precise mathematical proposition. This statement-only draft preserves that limitation and requires human review.
\subsection{Sources}
\begin{itemize}
\item Local UnsolvedMath registry, \texttt{data/unsolvedmath/problems.json}, permanent ID \texttt{30001975}, source code \texttt{catalogue entry}, title: \emph{Polanyi's Tacit Knowledge in Mathematical History}. The supplied record belongs to the Oberwolfach Reports Open Problems collection and gives proposed year 2012; no individual proposer is supplied.
\item Supplied report citation: \emph{Explicit Versus Tacit Knowledge in Mathematics} (2012), Oberwolfach Reports, pp.~131--210, DOI 10.4171/owr/2012/04. \url{https://doi.org/10.4171/owr/2012/04}. The background describes Oberwolfach Reports as documenting mathematical workshops and their open-problem discussions. No passage-specific page or independently verified attribution is supplied.
\item Supplied literature source, labeled ``original Oberwolfach report'': \url{https://doi.org/10.4171/owr/2012/04}. Literature checked at: \texttt{2026-08-21}. This is supplied dated curation derived from the OWR record, not an independently verified assessment. It labels the item \emph{open} and states that the extracted text poses a broad historical question about Polanyi and mathematical development, with no definitive answer that would justify a solved or disproved label. The supplied curation requests independent verification of time-sensitive status and cited claims.
\end{itemize}
\end{document}
